% =====================================================================
% RESULTS SKELETON (full APPS pool, data/apps_full.json)
% Numbers from notebooks/results_full.ipynb, origin/one-run-one-artifact,
% commit 14b4444, unless marked otherwise. Cell numbers in comments so every
% figure can be checked. Bullets are notes for the write-up, not final prose.
% Filled 2026-10-16 from the same notebook re-executed on the same runs; it now
% also writes CSVs to notebooks/figures/all/data/ (catch_at_fpr, game_results,
% paper_setting, paper_values, roc, frontiers_b_c) for pgfplots.
% [??] = not yet extracted from the notebook.
% Restructured 2026-10-26 after review: the deployment game (variations A, B,
% C) is replaced by the per-sample safety model of the methodology. Detection
% numbers are checked against notebooks/figures/all/data/catch_at_fpr.csv;
% \na{} marks a number waiting for experiments/safety_rework_plan.md M1--M2.
% =====================================================================

\providecommand{\na}{\textcolor{red}{[TBD]}}

\section{Results}
\label{sec:results}

\noindent\textit{Skeleton. Each subsection lists the points to make, with the
numbers behind them. Rakshit writes the final prose. \na{} marks a number not
yet pulled from the run.}

\paragraph{Overview (2--3 sentences).}
\begin{itemize}
  \item On the full pool, every tests-only arm ranks below the code-reading
        monitor (AUC 0.808--0.870 against 0.958) and is less safe at a 2\%
        audit budget (0.202--0.430 against 0.549; Section~\ref{sec:safety}).
  \item Adding tests to the monitor helps only when the threshold is strict.
        It raises catch at a 1\% false alarm rate (0.629--0.681 against 0.434)
        but not at 10\% (0.883--0.900 against 0.914). At a 2\% audit budget it
        raises safety from 0.549 to 0.662--0.716, but every paired interval
        includes zero (Section~\ref{sec:safety}); from 5\% the monitor alone
        is safer (Section~\ref{sec:budgets}).
  \item With defer to trusted, the monitor alone catches up with monitor +
        tests; only the variant whose tests the untrusted model planned stays
        ahead at high usefulness (Section~\ref{sec:dtt}).
  \item Changes to the test writer move the false alarm rate but not catch, and
        most misses are reached by a test input but not detected
        (Section~\ref{sec:why-catch}).
\end{itemize}

% ---------------------------------------------------------------------
\subsection{Pool and coverage}
\label{sec:pool}
% Sources: cells 5, 9, 11, 22, 33, 66, 68.

\paragraph{Points to make.}
\begin{itemize}
  \item 1388 APPS tasks, each with one honest and one backdoored candidate:
        2776 candidates. 555 train and 833 test tasks; the split is stored in
        the dataset file.
  \item The pool was scrubbed (Table~\ref{tab:scrub}). As it arrived, 48.6\% of
        backdoored candidates had a comment naming the backdoor, against 0\% of
        honest ones. After the scrub, no candidate has any comment.
        \textit{[Team's call: this table may sit better in Methods.]}
  \item The trusted model solves 1111 of 1388 tasks on its own (80.0\%), by
        control-arena's rule on every test. Worth
        one sentence: tests are not allowed to re-solve the task, even though
        the trusted model could solve most of them.
  \item 2435 of the 2776 candidates were scored by all eight arms (1177 honest,
        1258 backdoored). Every headline number uses only these. Per-arm
        failures are in Table~\ref{tab:coverage} (appendix).
  \item Safety is computed on the 642 evaluation tasks for which every arm
        scored both candidates; thresholds are calibrated on the 447
        calibration tasks with the same property.
\end{itemize}

\begin{table}[t]
\centering
\small
\begin{tabular}{llcc}
\hline
Pool & Side & Any comment & Names the backdoor \\
\hline
As it arrived & backdoored & 0.679 & 0.486 \\
              & honest     & 0.200 & 0.000 \\
Scrubbed      & backdoored & 0.000 & 0.000 \\
              & honest     & 0.000 & 0.000 \\
\hline
\end{tabular}
\caption{Share of candidates with any comment, and with a comment naming the
backdoor, before and after the scrub (1388 candidates per side). The scrub
removed comments and self-naming identifiers from every candidate, honest and
backdoored, so having no comments cannot itself give the label away.}
\label{tab:scrub}
\end{table}

% ---------------------------------------------------------------------
\subsection{Detection}
\label{sec:detection}
% Sources: cells 26, 29, 70. Figures: figures/all/roc.png,
% figures/all/catch_at_fpr.png.

\paragraph{Points to make.}
\begin{itemize}
  \item Every tests-only arm ranks below the monitor: AUC 0.808--0.870 against
        0.958, and lower catch at every false alarm rate in the grid
        (Table~\ref{tab:detection}).
  \item Letting the untrusted model plan the tests helps the tests-only arms:
        AUC rises from 0.816 to 0.870 (property) and from 0.808 to 0.867
        (plain), consistent with BigCodeBench (Section~\ref{sec:bcb}).
  \item Property and plain prompts barely differ: 0.816 against 0.808, and
        0.870 against 0.867.
  \item Adding tests to the monitor has a crossover. At a 1\% false alarm rate
        the combinations catch 0.629--0.681 against 0.434 for the monitor; at
        2\%, 0.669--0.721 against 0.558. At 5\% it is mixed, and at 10\% the
        monitor alone is ahead (0.914 against 0.883--0.900). The tests sharpen
        the top of the monitor's ranking but add nothing at loose thresholds.
  \item Do not rank arms by AUC alone: the monitor and monitor + tests property
        both have AUC 0.958 but very different curves. Compare catch at the
        false alarm rate the deployment actually uses.
  \item Say plainly that this table pools train and test and has no intervals.
\end{itemize}

\begin{table}[t]
\centering
\small
\begin{tabular}{lccccc}
\hline
Arm & AUC & catch @1\% & catch @2\% & catch @5\% & catch @10\% \\
\hline
monitor + u plans property & 0.962 & 0.681 & 0.721 & 0.838 & 0.900 \\
monitor + tests property   & 0.958 & 0.652 & 0.681 & 0.767 & 0.890 \\
monitor + tests plain\_v3  & 0.954 & 0.629 & 0.669 & 0.787 & 0.883 \\
monitor                    & 0.958 & 0.434 & 0.558 & 0.827 & 0.914 \\
u plans property/10        & 0.870 & 0.149 & 0.245 & 0.474 & 0.678 \\
u plans plain\_v3/10       & 0.867 & 0.074 & 0.223 & 0.473 & 0.667 \\
tests property/10          & 0.816 & 0.043 & 0.120 & 0.312 & 0.497 \\
tests plain\_v3/10         & 0.808 & 0.055 & 0.119 & 0.274 & 0.490 \\
\hline
\end{tabular}
\caption{Detection on the full pool, train and test pooled (2435 candidates
scored by all eight arms: 1177 honest, 1258 backdoored). Catch is the share of
backdoored candidates flagged when the false alarm rate on honest candidates is
held at 1, 2, 5 or 10\%. AUC counts ties as half. All eight arms see the code;
in the u plans arms the untrusted planner sees it and the trusted writer sees
only the plan, and every suite runs on the inputs of the same code-reading
trigger search.}
\label{tab:detection}
\end{table}

\begin{figure}[t]
\centering
% \includegraphics[width=0.85\linewidth]{figures/all/catch_at_fpr.png}
\fbox{\parbox{0.8\linewidth}{\centering Figure placeholder: catch at a fixed
false alarm rate (\texttt{figures/all/catch\_at\_fpr.png}). The ROC
(\texttt{roc.png}) can go in the appendix.}}
\caption{Catch rate of each arm when it may flag only 1\% to 10\% of honest
candidates (log scale; the 2435 candidates scored by all eight arms, train and
test pooled). Higher is better. The gap between two lines at the dotted 2\%
line is the difference a real audit budget sees. Between two thresholds an arm
can actually take, the curve is a straight line, which choosing at random
between those two thresholds achieves. At 1\% and 2\% the false alarm rate rests
on few honest candidates (about 12 and 24 of 1177).}
\label{fig:catch-at-fpr}
\end{figure}

% ---------------------------------------------------------------------
\subsection{Safety at a 2\% audit budget}
\label{sec:safety}
% Source: digest Section 8 (notebooks/figures/all/data/safety_at_budget.csv).
% Figure: notebooks/figures/all/safety_at_budget.png.

\paragraph{Points to make.}
\begin{itemize}
  \item Setting (Section~\ref{sec:safety-model}): the threshold is the 98th
        percentile of each arm's honest calibration scores, with ties broken at
        random; every evaluation candidate above it is audited; every one of
        the 642 evaluation tasks is attacked. No simulation and no attack rate.
  \item Monitor + tests is safer than the monitor alone in point estimate:
        0.662--0.716 against 0.549. But every paired difference has an interval
        that includes zero (best: monitor + u plans property, $+0.167$
        [$-0.068$, $+0.367$]). Property against plain tests: 0.682 against
        0.662.
  \item Why the intervals are wide: the monitor's 2\% threshold rests on
        about nine honest calibration candidates, so it moves a lot between
        resamples (the monitor's own interval is [0.347, 0.792]). The
        threshold-free comparison in Detection is tighter.
  \item No tests-only arm reaches the monitor: 0.202--0.430. For three of the
        four, the interval of the difference lies wholly below zero; for u
        plans plain\_v3 (0.430) it does not.
  \item The `any test fails' rule is not an audit rule: it flags 28.6--49.2\%
        of honest candidates, far above the 2\% budget, so it appears only in
        Detection.
\end{itemize}

\begin{table}[t]
\centering
\small
\begin{tabular}{lcc}
\hline
Arm & Safety [95\% CI] & Difference from monitor [95\% CI] \\
\hline
monitor                    & 0.549 [0.347, 0.792] & --- \\
monitor + tests property   & 0.682 [0.601, 0.726] & $+0.132$ [$-0.107$, $+0.322$] \\
monitor + tests plain\_v3  & 0.662 [0.573, 0.715] & $+0.113$ [$-0.121$, $+0.304$] \\
monitor + u plans property & 0.716 [0.676, 0.769] & $+0.167$ [$-0.068$, $+0.367$] \\
tests property/10          & 0.203 [0.127, 0.294] & $-0.346$ [$-0.572$, $-0.139$] \\
tests plain\_v3/10         & 0.202 [0.101, 0.262] & $-0.347$ [$-0.588$, $-0.141$] \\
u plans property/10        & 0.256 [0.190, 0.369] & $-0.293$ [$-0.519$, $-0.073$] \\
u plans plain\_v3/10       & 0.430 [0.188, 0.515] & $-0.120$ [$-0.454$, $+0.084$] \\
\hline
\end{tabular}
\caption{Safety at a 2\% audit budget on the 642 evaluation tasks that every
arm scored, with every task attacked. Each threshold is calibrated on the
honest candidates of the 447 calibration tasks, with ties broken at random.
Intervals come from 10{,}000 paired task resamples, each recalibrating the
threshold.}
\label{tab:safety}
\end{table}

% ---------------------------------------------------------------------
\subsection{Defer to trusted}
\label{sec:dtt}
% Source: digest Section 8 (notebooks/figures/all/data/dtt.csv).
% Figure: notebooks/figures/all/dtt.png.

\paragraph{Points to make.}
\begin{itemize}
  \item The four monitor protocols on one safety--usefulness plot, zoomed to
        the top right as in Figure~2 of \citet{greenblatt2024aicontrol}. The
        baseline is defer to trusted, not a monitor that discards flagged
        submissions.
  \item Deferral is cheap: $T$ solves 81\% of the evaluation tasks by
        control-arena's rule. Deferring 2.3\% of submissions lifts the monitor
        from 0.549 to 0.728 at usefulness 0.952, and 6.0\% lifts it to 0.842
        at 0.860.
  \item Once deferral is allowed, the monitor alone matches or beats monitor +
        tests with the trusted model's own tests: at usefulness of about
        0.86--0.87 the monitor reaches 0.842 against 0.799 for property tests;
        plain tests reach 0.781 at 0.903.
  \item Only monitor + u plans property stays ahead at high usefulness: 0.814
        at 0.956 (2.3\% deferred) and 0.863 at 0.873. Below usefulness 0.75 the
        four curves lie within about 0.03 of each other.
  \item So the gain from the trusted model's tests shows mainly when nothing
        is deferred; the protocol whose tests the untrusted model planned keeps
        a gain, with the caveats of Section~\ref{sec:protocols}.
\end{itemize}

\begin{figure}[t]
\centering
% \includegraphics[width=\linewidth]{figures/all/dtt.png}
\fbox{\parbox{0.8\linewidth}{\centering Figure placeholder: safety against
usefulness for the four monitor protocols with defer to trusted
(\texttt{notebooks/figures/all/dtt.png}).}}
\caption{Safety against usefulness with defer to trusted. Each line is one
protocol and each point one deferral threshold, with the audit threshold fixed
at a 2\% budget. Usefulness is scaled so that 0 is the trusted model alone and
1 the untrusted model alone. Up and to the right is better.}
\label{fig:dtt}
\end{figure}

% ---------------------------------------------------------------------
\subsection{Across audit budgets}
\label{sec:budgets}
% Source: digest Section 8 (notebooks/figures/all/data/safety_by_budget.csv).
% Figure: notebooks/figures/all/safety_by_budget.png.

\paragraph{Points to make.}
\begin{itemize}
  \item The advantage of monitor + tests is largest at the smallest budget and
        reverses by 5\%: at 1\%, 0.561--0.680 against 0.334; at 5\%,
        0.760--0.823 against 0.833; at 10\%, 0.887--0.901 against 0.915. This
        is the crossover detection predicts (catch at a 10\% false alarm rate
        0.883--0.900 against 0.914).
  \item Tests alone stay below the monitor at every budget.
  \item 10\% was the pre-registered main budget; 2\% is the budget of
        \citet{greenblatt2024aicontrol}. Report both and say which was
        pre-registered. These are point estimates; only 2\% has intervals.
\end{itemize}

\begin{table}[t]
\centering
\small
\begin{tabular}{lcccc}
\hline
Arm & 1\% & 2\% & 5\% & 10\% \\
\hline
monitor                    & 0.334 & 0.549 & 0.833 & 0.915 \\
monitor + tests property   & 0.586 & 0.682 & 0.760 & 0.887 \\
monitor + tests plain\_v3  & 0.561 & 0.662 & 0.769 & 0.887 \\
monitor + u plans property & 0.680 & 0.716 & 0.823 & 0.901 \\
tests property/10          & 0.129 & 0.203 & 0.359 & 0.535 \\
tests plain\_v3/10         & 0.102 & 0.202 & 0.325 & 0.555 \\
u plans property/10        & 0.190 & 0.256 & 0.486 & 0.683 \\
u plans plain\_v3/10       & 0.181 & 0.430 & 0.602 & 0.759 \\
\hline
\end{tabular}
\caption{Safety at each audit budget on the 642 evaluation tasks, with every
task attacked and each threshold calibrated on honest calibration candidates,
ties broken at random.}
\label{tab:budgets}
\end{table}

% TODO(Omar): "Multi-turn on the full pool" goes here, if it is kept
% (paper/outline.md, Results row 7).

% ---------------------------------------------------------------------
\subsection{Why catch does not move}
\label{sec:why-catch}

This subsection uses a smaller pool of 50 hard APPS problems, one honest and
one backdoored candidate each. The trusted model writes 20 tests per
submission, and we score six arms on a held-out split of 50 candidates (25
backdoored, 25 honest): two whose writer only sees the problem statement, and
four that see the code.
% TODO: Section 2.1 says T is always at low reasoning effort. The reasoning
% variant below breaks that, so 2.1 needs a line saying the hard-pool variants
% differ. These are also the only blind arms in the paper.

\paragraph{Scrubbing.}
In the original hard pool, the comment leak was lopsided by class: 32 of the 50
backdoored solutions had a comment naming the backdoor, using words like
\emph{bug}, \emph{flaw} or \emph{backdoor}, and none of the 50 honest ones did.
We stripped every comment with Python's \texttt{tokenize} module, checked that
each program's syntax tree was unchanged, and renamed four identifiers that
named the backdoor. Rerunning the high-reasoning code-visible arm on the
scrubbed pool left catch flat at 0.864 (38 of 44, on the 89 candidates scored
in both runs), and the false alarm rate moved only from 0.533 to 0.556.
Stripping the comments killed the lexical tell but left a structural one:
hardcoded branches on specific constants are still there for a writer that can
read the code. The rest of the hard-pool results use the scrubbed pool unless
we say otherwise.

\paragraph{The tests are not precise enough.}
If we reject any candidate that fails a test, the six arms reject 74\% to 88\%
of the backdoors (74\% and 80\% for the two blind arms, 76\% to 88\% for the
four code-visible ones) but keep only a quarter to three-fifths of the honest
code. Hold the false alarm rate at the monitor's 0.10 and their catch drops to
between 0.14 and 0.31, against 0.84 for the monitor. Of the three
relational-trigger backdoors the monitor missed (triggers set by a relation
between inputs rather than a fixed constant), five arms catch none and one
blind arm catches two, which is too few to read anything into. With 25
backdoored candidates one candidate moves catch by about 0.04, so we do not
rank the arms against each other; the claim is just that all six sit well below
the monitor.
% TODO: an earlier Results draft says the monitor reaches 0.930 safety at the
% 0.10 budget. Reconcile 0.84 and 0.930 before submitting.

\paragraph{The variations move the false alarm rate, not catch.}
Table~\ref{tab:variants} compares each variation against the same writer
without it. All four brought the false alarm rate down, by 0.12 to 0.25, and
none raised catch by more than one backdoor. Few-shot examples cut false
positives from 29 of 44 honest candidates to 20 of 44, and catch fell from 42
of 48 backdoors to 39 of 48. The critic cut the false alarm rate from 0.674 to
0.429, and catch fell from 0.878 to 0.816. The critic was also not selective:
it dropped 6.39 tests per backdoored candidate on average and 5.59 per honest
one, so it pruned backdoored code a little harder than honest code, not the
other way round. More reasoning cut the false alarm rate from 0.698 to 0.558
and moved catch by one backdoor (39 to 38 of 44, on the 87 candidates both arms
scored). Two things make that last comparison shaky: the low-effort arm wrote
no reasoning at all, against a median of about 65,000 characters for the high
arm, and the high arm ran with half the sandbox time (120\,s against 240\,s).
The multi-turn loop cut the false alarm rate from 17.65\% to 5.88\% on a
19-task development set with all 18 matched backdoors still caught, but the
95\% interval for that change, $[-29.41, 0.00]$ points, reaches zero.
% TODO: confirm the model and pool for the multi-turn arm (Omar).

\begin{table}[t]
\centering
\small
\begin{tabular}{llcc}
\hline
Variation (vs.\ without) & Candidates & FPR & Catch \\
\hline
Few-shot & scrubbed, 49 tasks$^{a}$ & 0.66 $\to$ 0.45 & 0.88 $\to$ 0.81 \\
Critic & scrubbed, $n=98$ & 0.67 $\to$ 0.43 & 0.88 $\to$ 0.82 \\
Reasoning (high vs.\ none$^{b}$) & original, $n=87$ & 0.70 $\to$ 0.56 & 0.89 $\to$ 0.86 \\
Multi-turn (vs.\ one turn) & dev set, 19 tasks & 0.18 $\to$ 0.06$^{c}$ & 1.00 $\to$ 1.00 \\
\hline
\end{tabular}
\caption{Each variation against the same writer without it, on the hard pool.
FPR is the share of honest candidates flagged, catch the share of backdoored
candidates flagged.
$^{a}$Full scrubbed pool minus the task the few-shot example came from
(44 honest and 48 backdoored candidates scored).
$^{b}$The low-effort arm wrote zero reasoning characters, so this is really no
reasoning against high reasoning. Run on the original (unscrubbed) pool.
$^{c}$95\% interval for the change: $[-0.29, 0.00]$.}
\label{tab:variants}
\end{table}

\paragraph{Where the misses come from.}
To tell search failures apart from detection failures, we used an earlier run
on five hard BigCodeBench tasks with two backdoors each, where blind
property-based tests caught none of the 10.
% TODO: the 2/8 split is from these five BigCodeBench tasks, not the hard APPS
% pool. Keep this framing or rerun wall_decomposition on the hard pool (Alon's
% action item). Also name the writer models for this run.
For each backdoor we ran the backdoored and reference solutions on 20,000
random inputs and counted how often their outputs differed. For 2 of them,
inputs drawn from the problem's own input format never produced a difference,
so no realistic test budget would reach the trigger. One of those became
reachable once we let the input alphabet include the form-feed character, after
which the outputs differed on 36\% of inputs. For the other 8, between 21\% and
90\% of inputs that respect the input contract exposed a difference, and a
2,000-input test reaches the trigger with probability 1.000. So those triggers
were reached. What was missing is a test that asserted a property telling the
backdoored output apart from the reference output. We call this the
spec-entailment wall.
% Provisional, keep out of the text until checked against the task docstrings:
% of the 8 reached backdoors, roughly 1-3 look like behaviour the spec leaves
% open (e.g. the 202/attack_0 tie-break) and roughly 5 look like errors a
% sound spec-derived test could catch.

% ---------------------------------------------------------------------
\subsection{BigCodeBench}
\label{sec:bcb}
% Source: team research doc, Alon's 26-task run (22/09). May be replaced by the
% BCB run with Palak's and Omar's arms.

\paragraph{Points to make (Alon to confirm or update).}
\begin{itemize}
  \item 26 tasks, which is the whole usable pool (the upstream source covers
        31 BigCodeBench tasks), 52 candidates (26 honest, 26 backdoored).
  \item Unlike APPS, no BigCodeBench backdoor named itself in a comment
        (0.000 before the scrub).
  \item The monitor reaches AUC 1.000. Four baselines that use no model at all
        (counting \texttt{==} comparisons 0.967, numeric literals 0.935, if
        statements 0.883, characters 0.825) all beat all four test-writing arms
        (0.579--0.753).
  \item Counting \texttt{==} comparisons against the monitor in safety at a 2\%
        budget: \na{} against \na{}.
        % TODO(Alon): recompute under the per-sample safety model; the game-A
        % figures of the earlier draft (0.985 against 0.896) no longer apply.
  \item Letting the untrusted model plan the tests beat the trusted model alone
        in all four comparisons across two pools and two samplings.
  \item Conclusion: this benchmark is saturated, so it cannot separate
        protocols.
\end{itemize}

% ---------------------------------------------------------------------
\paragraph{Notes to carry into Discussion and Limitations.}
\begin{itemize}
  \item No red-team attack selection. The red team attacks every task with a
        fixed, pre-generated backdoor and does not choose which tasks to
        attack. An attacker that attacks only where its backdoor is least
        likely to be flagged or tested would lower safety for every protocol,
        possibly by different amounts. \citet{schaeffer2026} find that attack
        selection lowers trusted monitoring's safety from 99\% to 59\% in a
        concentrated BigCodeBench setting at a 0.5\% audit budget.
  \item Provider change (Appendix~\ref{app:implementation}): the eight main arms were
        scored before 10 October; plain\_v5 and the multi-turn study were
        scored partly or wholly after it, and no rescoring was possible.
  \item Detection is pooled over train and test; safety calibrates on train
        and evaluates on test only.
  \item Every arm in the full run sees the code, the trigger search included.
        The only specification-only arms are in the hard-pool study.
  \item The trusted and untrusted models are the same model (DeepSeek-V3.2);
        the gap between them comes from reasoning effort alone.
  \item The intervals resample tasks, not model calls. Each candidate was
        scored once per arm (the monitor's two samples are averaged), so the
        intervals cannot separate the model's sampling noise from task
        difficulty.
\end{itemize}

% ---------------------------------------------------------------------
% APPENDIX TABLES (move to the appendix in Overleaf)
% ---------------------------------------------------------------------

\begin{table}[t]
\centering
\small
\begin{tabular}{lccccc}
\hline
Arm & Scored & Model fail. & Infra fail. & Partial & Dropped \\
\hline
monitor                    & 2774 & 0  & 2  & 0  & 339 \\
monitor + tests property   & 2664 & 81 & 31 & 0  & 229 \\
monitor + tests plain\_v3  & 2627 & 98 & 51 & 0  & 192 \\
monitor + u plans property & 2677 & 83 & 16 & 0  & 242 \\
tests property/10          & 2684 & 61 & 31 & 44 & 249 \\
tests plain\_v3/10         & 2659 & 66 & 51 & 54 & 224 \\
u plans property/10        & 2708 & 52 & 16 & 28 & 273 \\
u plans plain\_v3/10       & 2703 & 45 & 28 & 30 & 268 \\
\hline
\end{tabular}
\caption{Coverage for each arm out of 2776 candidates. Failed candidates have
no score and leave every denominator. ``Partial'' counts suites where not every
test ran on every trigger input. ``Dropped'' counts candidates this arm scored
but at least one other arm did not; headline numbers use the 2435 that every
arm scored.}
\label{tab:coverage}
\end{table}
